\abstract{\textbf{Motivation:} Peptide-MHC-I binding-prediction benchmarks can manufacture false conclusions when confounders are uncontrolled.\
\textbf{Results:} We benchmarked six architectures on 5,190 MHCflurry-labelled HLA-A*02:01 9-mers and re-analysed three apparent findings under controlled conditions. A convolutional network was highest (91.9\% vs.\ 90.6\% deep feed-forward); random labels fell to chance, and a claimed ``label quality dominates architecture'' result dissolved once the degenerate null was removed. A zero-parameter anchor rule matched the trained model on canonical peptides, while non-canonical binders plateaued at $\sim$0.65 sensitivity across four predictors --- the true frontier. Extending to five HLA-I alleles and a 15-allele panel, a rebalanced pan model matched per-allele specialists on four of five alleles (0.819 pooled non-canonical sensitivity). A temporal holdout preserved post-2022 recognition; 6,524 experimental affinities (AUC 0.942) confirmed label self-consistency.\
\textbf{Availability and Implementation:} Code, models, and data at \href{https://github.com/woodhaha/peptide-mhc-benchmark}{github.com/woodhaha/peptide-mhc-benchmark}.\
\textbf{Contact:} \href{zhouzhuha@wmu.edu.cn}{zhouzhuha@wmu.edu.cn}; \href{wzxxy@wmu.edu.cn}{wzxxy@wmu.edu.cn}\
\textbf{Supplementary Information:} Supplementary data are available at \textit{Bioinformatics} online.}
